\documentclass{article}

\usepackage{amsmath, amssymb}
\usepackage{graphicx}
\usepackage{booktabs}
\usepackage{hyperref}
\usepackage{url}

\title{
HASIP: Hardware-Aware Structured Pruning for Energy-Efficient Neural Network Inference
}

\author{
Valentino Lanzarini \\
\texttt{Experimental Research Framework}
}

\date{April 2026}

\begin{document}

\maketitle

\begin{abstract}
We present HASIP (Hardware-Aware Structured Inference Pruning), a framework for analyzing the relationship between structured pruning, hardware execution behavior, and energy consumption in neural network inference. Through controlled experiments on a fully connected neural network deployed on GPU (NVIDIA Tesla T4) and CPU backends via PyTorch and ONNX Runtime, we evaluate how structured sparsity affects latency, throughput, and estimated energy consumption. Results show that while pruning consistently reduces parameter count (up to 91.4\%) and FLOPs (up to 91.5\%), latency scaling is non-linear and hardware-dependent. GPU performance remains dominated by kernel launch overhead in small-FLOP regimes, while energy consumption decreases monotonically up to 97.8\%. These findings highlight that model compression primarily optimizes energy efficiency rather than latency in low-compute inference regimes.
\end{abstract}

\section{Introduction}

Neural network compression techniques such as pruning and quantization are widely used to reduce computational cost. However, their effect on real hardware performance is non-trivial, especially in small-scale inference regimes where hardware overhead dominates computation.

In this work, we study the interaction between structured pruning and hardware execution behavior. Unlike prior works focusing solely on model accuracy or FLOPs reduction, we analyze:

\begin{itemize}
    \item Latency across CPU and GPU backends
    \item Energy consumption under analytical modeling
    \item Hardware utilization efficiency under sparsity
\end{itemize}

We introduce HASIP, a structured experimental pipeline for evaluating these effects.

\section{Related Work}

Model compression has been extensively studied through pruning \cite{han2015deep}, quantization \cite{jacob2018quantization}, and compiler-level optimization \cite{chen2018tvm}. Energy-aware machine learning has also been explored in Green AI literature \cite{schwartz2019green}.

However, most prior work evaluates either:
\begin{itemize}
    \item model-level efficiency (FLOPs, parameters), or
    \item system-level inference latency
\end{itemize}

but not their interaction under structured sparsity.

\section{Methodology}

\subsection{Model Architecture}

We evaluate a fully connected neural network:

\[
100 \rightarrow 512 \rightarrow 256 \rightarrow 10
\]

with ReLU activations.

\subsection{Structured Pruning}

We apply layer-wise structured pruning based on neuron importance:

\[
s_i = \frac{1}{n} \sum_j |W_{ij}|
\]

Neurons are removed according to percentile thresholding.

\subsection{Metrics}

We measure:

\textbf{Latency:}
\[
L = \frac{1}{N} \sum t_i
\]

\textbf{Speedup:}
\[
S = \frac{L_{base}}{L_{pruned}}
\]

\textbf{Energy (estimated):}
\[
E = \alpha \cdot FLOPs + \beta \cdot Memory
\]

\section{Experimental Setup}

\textbf{Hardware:}
\begin{itemize}
    \item GPU: NVIDIA Tesla T4 (14.6 GB VRAM)
    \item Framework: PyTorch 2.10 + CUDA 12.8
    \item Backend: ONNX Runtime (CPU + GPU)
\end{itemize}

\textbf{Model:}
\begin{itemize}
    \item Parameters: 185,610 (baseline)
    \item FLOPs: 369,664
    \item Batch size: 1024
\end{itemize}

\textbf{Pruning levels:}
\[
0\%, 20\%, 40\%, 60\%, 80\%
\]

\section{Results}

\begin{table}[h]
\centering
\caption{Experimental results on Tesla T4}
\begin{tabular}{lccccc}
\toprule
Model & Params & FLOPs & Latency (ms) & ONNX GPU (ms) & Energy (µJ) \\
\midrule
Base & 185k & 369k & 0.247 & 0.368 & 149.9 \\
20\% & 127k & 252k & 0.136 & 0.373 & 102.5 \\
40\% & 79k  & 157k & 0.101 & 0.347 & 64.1 \\
60\% & 42k  & 84k  & 0.178 & 0.292 & 34.2 \\
80\% & 15k  & 31k  & 0.101 & 0.252 & 12.9 \\
\bottomrule
\end{tabular}
\end{table}

\section{Key Findings}

\subsection{1. Non-linear latency scaling}
Latency reduction is not proportional to pruning rate. Best efficiency occurs between 20–40\% sparsity.

\subsection{2. GPU inefficiency in small models}
For low FLOP regimes, GPU execution is dominated by kernel launch overhead rather than computation.

\subsection{3. CPU competitiveness}
ONNX CPU runtime becomes competitive due to lower overhead and better utilization in small models.

\subsection{4. Energy reduction is monotonic}
Energy decreases consistently with pruning, reaching up to 97.8\% reduction.

\section{Discussion}

These results suggest a decoupling between:

\begin{itemize}
    \item computational complexity (FLOPs)
    \item runtime latency
    \item energy consumption
\end{itemize}

In small neural networks, hardware overhead dominates performance more than arithmetic operations.

\section{Limitations}

\begin{itemize}
    \item Results are based on a small MLP architecture
    \item Energy is estimated, not physically measured
    \item No TensorRT or low-level kernel optimization used
\end{itemize}

\section{Conclusion}

We introduced HASIP, a hardware-aware framework for evaluating structured pruning effects on inference efficiency. Our experiments show that:

\begin{itemize}
    \item pruning strongly reduces energy consumption
    \item latency improvements are hardware-dependent
    \item GPU acceleration is not guaranteed in low-compute regimes
\end{itemize}

These findings highlight the importance of considering hardware execution dynamics in model compression research.

\begin{thebibliography}{9}

\bibitem{han2015deep}
S. Han et al.,
"Deep Compression",
ICLR 2016.

\bibitem{jacob2018quantization}
B. Jacob et al.,
"Quantization and Training of Neural Networks",
CVPR 2018.

\bibitem{chen2018tvm}
T. Chen et al.,
"TVM: An Automated End-to-End Optimizing Compiler",
OSDI 2018.

\bibitem{schwartz2019green}
R. Schwartz et al.,
"Green AI",
arXiv 2019.

\end{thebibliography}

\end{document}
